Sources are collected in a literature database. The word
"database"' may give rise to some misconceptions. For \LaTeX\ one understands
a simple text file in the so-called \hologo{BibTeX} format.
Listing~\ref{lst:latex:bibliography:bibtexentry} shows an example entry for this
manual.

\lstinputlisting[language=BibTeX,style=block,float,caption={Ein \hologo{BibTeX}-Eintrag},label={lst:latex:bibliography:bibtexentry}]{code/latex_bibliography_bibtexentry.bib}

An entry consists of three parts.
%
\begin{compactenum}
\item the type, here \texttt{@manual}
\item a key, here \texttt{urz\_vorlage\_manual}
\item several attributes, which describe the locations in more detail.
\end{compactenum}

Of course, you can maintain the file manually with a text editor. However,
there is specialised software that allows the editing of a \hologo{BibTeX}
database. Examples are:
%
\begin{compactitem}
\item K\hologo{BibTeX} (\url{https://userbase.kde.org/KBibTeX})
\item Zotero (\url{https://www.zotero.org})
\item JabRef (\url{https://www.jabref.org})
\end{compactitem}
%
Sometimes these programmes save the literature database in a separate format
and export the format and export the references after processing only in the
\hologo{BibTeX} format.

In the rarest cases it is necessary to create a \hologo{BibTeX}-database
manually. Publishers and libraries usually have a corresponding export function
in their catalogues, which makes it possible to create prefabricated \hologo{BibTeX}
entries. However, subsequent adaptations are generally adjustments are
necessary if umlauts or special characters in the author names are not encoded
correctly or data fields are missing. 
